\section{Querschnittliche Konzepte}

Die folgenden technischen Regeln gelten serviceübergreifend.

\subsection{Sicherheit}

Der Auth Service stellt JWTs mit einer Gültigkeit von 60 Minuten aus.
Das API Gateway validiert die Tokens und übergibt die Benutzer-ID über den Header \texttt{X-User-Id}.
Der interne Aufruf des User Service an den Auth Service wird zusätzlich durch den Header \texttt{X-Internal-Key} autorisiert.
\subsection{Kommunikation}

Synchrone Serviceaufrufe erfolgen über REST.
Bestellereignisse werden über MQTT mit Zustellbestätigung übertragen.

\subsection{Fehlerbehandlung}

Die Spring-basierten Services verwenden ein einheitliches Fehlerformat.
Validierungsfehler liefern \texttt{400}, fehlende Ressourcen \texttt{404}, Konflikte \texttt{409} und
unerwartete Fehler \texttt{500}.

\subsection{Datenhaltung}

Jeder zustandsbehaftete Service besitzt eine eigene PostgreSQL-Datenbank.
Nur der jeweilige Service darf direkt auf seine Datenbank zugreifen.
Serviceübergreifende Beziehungen werden über IDs statt über
datenbankübergreifende Fremdschlüssel abgebildet.

\subsection{Serviceübergreifende Konsistenz}

Serviceübergreifende Datenbanktransaktionen werden nicht verwendet.
Schlägt die Änderung eines Produkts während des Kaufvorgangs fehl, setzt der
Order Service bereits vorgenommene Änderungen durch kompensierende Aufrufe
zurück.

\subsection{Skalierung}

Die Instanzen des Product Service speichern keine Sitzungsdaten oder lokalen
Dateien. Dadurch können Anfragen über Nginx im Round-Robin-Verfahren auf drei
Replikate verteilt werden.

Alle Replikate greifen auf dieselbe \texttt{product-db} zu. Die Datenbank
kann daher bei hoher Last zum Engpass werden. Eine automatische,
lastabhängige Skalierung ist nicht umgesetzt; die Anzahl der Replikate wird
beim Start festgelegt.